\section{Introduction}

A periodic orbit can carry arithmetic data that are invisible after its
parameter or action coordinate has been eliminated.  The HCS-C33
construction made this phenomenon exact for a period-five collision in an
area-preserving H\'enon family: two distinct chronological orbits have the
same generating action, while the product of their two Hill determinants
descends to a nontrivial square class over the collision field
\citep{hcs2026c33note,hcs2026c33certificate}.  That theorem produced one
quadratic extension.  It did not decide what happens when all nine Galois
conjugate collision parameters are placed in a common splitting field.

The missing question is a rank problem rather than another orbit search.
Let \(P_9\) be the degree-nine collision polynomial, let
\(K=\Q[A]/(P_9)\), and let \(L\) be its splitting field.  C33 proves
\(\Gal(L/\Q)=S_9\).  If \(\beta\in K^\times\) is the symmetric
two-branch Hill product and \(\beta_1,\ldots,\beta_9\) are its conjugates
in \(L\), define
\[
 V=\langle[\beta_1],\ldots,[\beta_9]\rangle
 \subset L^\times/L^{\times2}.
\]
C33 proves only \(V\ne0\).  Hidden products of conjugates could still make
\(\dim_{\F_2}V\) any of several smaller values.  Irreducibility of one
degree-eighteen polynomial would not, by itself, exclude those products.

This note closes the rank problem at its maximal value.

\begin{theorem}[Main theorem]\label{thm:intro-main}
For the exact C33 period-five Maxwell--Hill object,
\[
 \dim_{\F_2}\langle[\beta_1],\ldots,[\beta_9]\rangle=9.
\]
Consequently, with
\(M=L(\sqrt{\beta_1},\ldots,\sqrt{\beta_9})\),
\[
 \Gal(M/\Q)\cong C_2\wrp S_9,
 \qquad |\Gal(M/\Q)|=2^9 9!=185{,}794{,}560.
\]
\end{theorem}

The proof has three short layers.  First, a translation
\(A=1802+T\) exposes a lower \(19\)-adic Newton edge from \((0,5)\) to
\((2,0)\).  The two roots in that cluster both give odd, integer-normalized
valuation \(5\) to \(\beta\); the other seven Hill values are units.  This
produces the parity vector \(e_1+e_2\).  Second, conjugation by \(S_9\)
produces every vector \(e_i+e_j\), forcing a square relation to have all
coordinates equal.  Third, the all-ones relation would make the rational
norm of \(\beta\) a square in \(L\).  Its square-free class is
\(3\cdot13\cdot19\cdot41\cdot59\), whereas the sign quadratic subfield has
class \(13\cdot19\cdot41\cdot59\).  Their quotient has class \(3\), so the
last relation is impossible.

The concrete contributions are therefore:

\begin{enumerate}[leftmargin=*,itemsep=2pt]
\item an exact local weight-two parity certificate for the intrinsic
      H\'enon Hill product;
\item a proof that the resulting nine-dimensional Kummer module has no
      hidden relation in the common \(S_9\) splitting field;
\item an order-sharp upgrade of the standard quadratic wreath embedding to
      equality; and
\item a reproducible exact certificate that separates these statements
      from all-period zeta and Hilbert--P\'olya claims.
\end{enumerate}

The result is a substantial finite arithmetic structure, but its scope is
narrow.  No prime is fitted to a Riemann-zero target: \(19\) is forced by
the exact norm/discriminant ledger and used only as a local proof place.
No Euler product, analytic continuation, functional equation, critical-line
law, or self-adjoint operator follows from the finite Galois group.
